

\illustration{0.3}{cow.jpeg}{
   A cow.
   CC BY-SA 3.0 by Kim Hansen on \href{https://commons.wikimedia.org/wiki/File:Cow_(Fleckvieh_breed)_Oeschinensee_Slaunger_2009-07-07_(cropped).jpg}{Wikimedia~Commons}
}





Farmer Kiki\ has $42$ cows.
Their field can be thought of as a grid where each cell is a square with a side length of one metre.
According to strict regulations of the German Cow Perimeter Conglomerate, the cows must graze on a fenced area that is exactly $a$ square metres large.

Unfortunately, Farmer Kiki's supplier messed up their order of fence pieces.
Instead of mostly straight fence pieces and just a few corner fence pieces, they only received an enormous amount of corner fence pieces.
There are even more pieces than would fit on the field.
One corner fence looks like this:

\texttt{.... \\
.O\#. \\
.\#.. \\
....}

Here `\texttt{O}' denotes the center of the fence, `\texttt{\#}' denotes the sides of the fence, and `\texttt{.}' denotes empty space.

Farmer Kiki~can rotate and move each corner fence grid-aligned.
How may they place them such that there is a side-connected component of exactly $a$ empty cells that is surrounded by fence pieces?
More exactly, for each interior cell of that component all four orthogonal neighbours have to either also be an interior empty cell or part of a fence piece.

Farmer Kiki's field is $100 \times 100$ cells large and they cannot place any fence piece outside this area.

\begin{figure}[h]
    \centering
    \includegraphics[height=4cm]{figure1}
    \includegraphics[height=4cm]{figure2}
    \includegraphics[height=4cm]{figure3}
    \caption{Two valid and one invalid construction for the first sample with fenced area~$4$. The left construction represents the given sample answer. On the right, the four interior cells are not side-connected.}
\end{figure}

